\chapter{Reference map, Lean handoffs, and quarantined routes}
\label{route-app-technical}

This appendix is part of the authoritative specification.  It records where
each mathematical assertion comes from, what exact Lean object receives it,
and which nearby constructions are deliberately outside the route.

\section{Reference map}
\label{route-app-reference-map}

\begin{center}
\begin{tabular}{p{0.08\textwidth}p{0.32\textwidth}p{0.50\textwidth}}
\textbf{Step} & \textbf{Precise anchor} & \textbf{Specialization used here}\\ \hline
R1 & Kleiman, \emph{The Picard Scheme}, Section 3, Definition 3.6,
Theorem 3.7, and Exercise 3.8. & Relative effective divisors commute with
base change; for a projective flat curve their functor is represented by an
open Hilbert subscheme, and fixed degree is open-and-closed and of finite
type.  Lean proves the analogous statement only for its exact widened,
locally certified degree-\(g\) functor.\\
R2 & Kleiman, Section 3, Theorem 3.13; Milne, \emph{Abelian Varieties},
Chapter III, Theorem 5.1(a) and Lemma 5.2(b), transcription page 0107. & An
Abel fibre is a complete linear system.  Under \(H^1=0\) and rank-one
pushforward, its projective bundle is the base, giving one effective divisor.
The checked familywise evaluator and uniqueness theorem provide the actual
Lean specialization.  Milne III.5 additionally assumes complete nonsingular,
positive genus, and a rational point; those are not AJCR route hypotheses.\\
R3 & Stacks, Varieties Lemma 33.25.6, Tag 056U; Milne III, Lemma 5.2(b). & Tag
056U gives density of closed points with residue field finite separable over
the base.  Its use here is a derived specialization: restrict to the required
nonempty smooth open, choose such a point, and then use separable closedness to
split the finite separable residue extension.  This is the rational-point
input inside translated-drop coverage, while Milne supplies the
classical degree-\(g\), rank-one model subject to the same III.5 caveat.\\
\end{tabular}
\end{center}

\begin{center}
\begin{tabular}{p{0.08\textwidth}p{0.32\textwidth}p{0.50\textwidth}}
\textbf{Step} & \textbf{Precise anchor} & \textbf{Specialization used here}\\ \hline
R4 & Stacks, Properties of Schemes Lemma 26.15.4, Tag 01JJ; Morphisms Lemma
29.22.9, Tag 01TX; Kleiman, Theorem 4.8(1). & A sheaf covered by representable
open subfunctors is represented; over the Noetherian base \(\Spec k\), local
finite type implies local finite presentation.  Kleiman's full Picard theorem
is context, not an imported shortcut.\\
R5 & Stacks, Schemes Lemmas 26.14.1--2, Tags 01JB--01JC; Limits Lemma
32.10.1, Tag 01ZM; Limits Lemmas 32.8.12 and 32.8.11, Tags 0EUU and 081E. &
Descend the finite-presentation charts, maps, and equalities to a common stage;
Tag 01ZM is applied to an inverse system with affine transition maps and qcqs
stages, while Tags 0EUU and 081E require the relevant stage maps to be locally
of finite presentation before openness or invertibility is reflected.  Scheme
gluing and its mapping property yield the global carrier comparison.  None of
these results proves Picard representability over the stage.\\
R6 & Stacks, Fields Lemmas \emph{normal-closure-inside-normal}, Tag 0BMG, and
\emph{normal-closure-galois}, Tag 0EXM; Picard Schemes of Curves Section
44.4, Tag 0B9K; Descent Remark 35.35.12, Tag 040L. & Take the finite normal
closure of all coefficients inside the fixed normal extension \(\Omega/K\)
before finalizing \(P\): Tag 0BMG supplies containment and Tag 0EXM supplies the
Galois conclusion for the finite separable input.  Then descend the class and
prove the all-tests Yoneda equivalence on \(P.\mathrm{gluedOver}\) using
arbitrary-test Picard reflection and descent of the representing morphism
with its double-overlap compatibility.  Classes are constructed in AJCR's
big-etale sheaf itself; Tag 0B9K describes the fppf Picard sheaf and is
background rather than an identification of the two functors.\\
\end{tabular}
\end{center}

\begin{center}
\begin{tabular}{p{0.08\textwidth}p{0.32\textwidth}p{0.50\textwidth}}
\textbf{Step} & \textbf{Precise anchor} & \textbf{Specialization used here}\\ \hline
R7 & Stacks, Groupoid Schemes Lemma 39.8.7, Tag 0BF7; Properties of Schemes
Lemma 28.30.5, Tag 01ZY; Algebra Lemma 10.57.6, Tag 00JS. & Representation
and finite type make the literal \(J_L=P.\mathrm{gluedOver}\) an algebraic
group.  Tag 0BF7 gives quasi-projectivity over arbitrary \(L\), and
homogeneous prime avoidance gives affine neighbourhoods of finite orbits.
The field-general group theorem remains an unformalized source prerequisite.\\
R8 & Stacks, Descent Section 35.38, Tag 0246, and Remark 35.35.12, Tag 040L. &
Galois-stable quasi-affine charts descend effectively, and compatible
morphisms descend uniquely.  The checked finite-Galois quotient theorem
specializes this to the represented degree-zero Picard functor.\\
R9 & Milne III, Proposition 6.1, transcription page 0110; Kleiman, Section 5,
especially Lemma 5.1 and Corollaries 5.13--5.14. & Milne proves the pointed
Albanese property for positive genus; the challenge also requires genus zero.
Kleiman supplies the classical group-scheme geometry.  The datum-level Lean
consumers retain every extra hypothesis explicitly.\\
\end{tabular}
\end{center}

\section{Exact Lean handoffs}
\label{route-app-lean-handoff}

\begin{center}
\begin{tabular}{p{0.08\textwidth}p{0.46\textwidth}p{0.36\textwidth}}
\textbf{Step} & \textbf{Public handoff and exact carrier} & \textbf{State}\\ \hline
R1 & \texttt{divFunctorAff\_genus\_representableBy} on
\texttt{divRepAffGenusScheme C}, representing exactly
\texttt{divFunctorAff C (genus C)}. & Proved.\\
R2 & \texttt{canonicalRankOneAbelIso} between the two sigma-extended loci;
\texttt{canonicalRankOneAbelSliceIso} slicewise; and
\texttt{picRankOneOpenSigmaIso} from the represented open
\texttt{divRankOneOpenOver}. & Proved.\\
R3 & \texttt{picRankOneTranslatedChart\_pointwiseCoverage} for the exact
translator subtype \texttt{PicRankOneTranslatorIndex}. & Proved.\\
R4 & \texttt{pic0\_sepClosed\_representableBy C} with carrier
\texttt{.1} and certificate \texttt{.2}; the selected universal class is
\texttt{pic0SepClosedUniversalClass C}. & Proved.\\
R5 & For arbitrary legacy
\texttt{P0 : Pic0FiniteStageGluePackage C\_Omega K}, carrier
\texttt{P0.gluedOver}; comparison \texttt{P0.finiteStageBaseChangeIso} plus
its structure-map equation. & The global comparison remains unverified on
current source.  Its target is only a carrier comparison; the package supplies
no separability or Galois certificate.\\
\end{tabular}
\end{center}

\begin{center}
\begin{tabular}{p{0.08\textwidth}p{0.46\textwidth}p{0.36\textwidth}}
\textbf{Step} & \textbf{Public handoff and exact carrier} & \textbf{State}\\ \hline
R6 & Jointly choose legacy \texttt{P} with \texttt{L = P.N.1} finite Galois,
then construct
\texttt{xiL : pic0Subgroup ((baseChange K L).obj C) P.gluedOver}, (UC), and
the exact \texttt{RepresentableBy} certificate. & Unresolved; the class is the
first unmet mathematical producer.\\
R7 & Prove quasi-projectivity and \texttt{FiniteInAffine} on the represented
finite-type group \texttt{P.gluedOver}, then construct
\texttt{rho.OrbitsInAffineOpen} for the canonical semilinear action
\texttt{rho} induced by \texttt{repL}. & Group quasi-projectivity is an independent unformalized
theorem; its exact-carrier application additionally requires R6.
Projectivity wrappers are stronger conditional consumers only.\\
R8 & \texttt{pic0RepresentableBy\_finiteGaloisDescent C repL}; carrier
\texttt{StableAffineOpen.gluedQuotientOver ...}. & Proved conditionally on
\texttt{repL} and orbit-affineness.\\
R9 & \texttt{picRepDatum\_finiteGaloisDescent}, then
\texttt{PicRepDatum.toJacobianData}; all shared fields preserve the quotient
carrier by reflexivity. & Proved conditionally; original-field instantiation
blocked by R6--R7.\\
P8 & \texttt{JacobianData.ofCurve}, \texttt{pullbackHom},
\texttt{baseChange}, the tangent, geometry, Abel-base-change, and conditional
Albanese families, then the declarations in \texttt{Challenge.lean}. & Datum
consumers proved with displayed inputs; their first unmet producers and the
still-sorry-bearing headline assembly are recorded below.\\
\end{tabular}
\end{center}

\section{Quarantine}

\begin{remark}[Rejected and non-normative branches]
  \label{route-app-quarantine}
  \dcref{custom}

  The following material is not an alternate choice left to a formalizer.
  \begin{enumerate}
  \item The unrestricted high-degree or admissible Abel map, its quotient by
    complete linear systems, and the symmetric-power/Weil or Albanese-first
    constructions are rejected as construction routes.  Milne's
    \(n>2g\) argument cannot justify the degree-\(g\) rank-one step.  The fixed
    admissible integer may be reused only in the downstream Phase 8 producer
    of an \texttt{AbelSourceData}; it does not replace R1--R8 or represent
    \(\operatorname{Pic}^0\).
  \item \texttt{DivFamZar}, the conditional F5/F6 faithfully-flat divisor
    classifiers, and the admissible-degree representer are implementation
    layers or quasi-compactness inputs only.  They do not replace R1--R2.
  \item A \texttt{GlueData}, \texttt{GluePackage}, overlap isomorphism, or
    \texttt{pic0RepresentableBy\_finiteStageBaseChange} is not a
    representability theorem over the finite stage.  Only (UC) followed by
    (Y) may cross R6.
  \item The canonical \texttt{Pic0FiniteStageGluePackage} is the sole
    finite-stage carrier on R5--R7.  Orbit and descent consumers use it
    directly; no context wrapper, stable wrapper, or carrier transport is
    part of the selected route.
  \item The Lean functor \texttt{pic0TypeFunctor} means the degree-zero
    etale-sheafified Picard functor.  No theorem currently identifies it with
    Kleiman's connected identity component, so those phrases are not
    interchangeable in Lean attachments or completion claims.
  \item The R7 immersion is obtained by applying algebraic-group
    quasi-projectivity to the represented final carrier.  Projectivity,
    quasi-projectivity, and finite-in-affine are distinct properties;
    projectivity is not an additional prerequisite for R8.  Kleiman Theorem 4.8,
    arbitrary object-isomorphism transport,
    post-base-change representability, and filtered-colimit preservation are
    comparison or consumer results.  None is an imported proof of R6 or R8.
  \item The conditional constructors
    \texttt{JacobianData.ofRepresentableBy}, \texttt{ofCharts},
    \texttt{ofChartsOfCompactSpace}, \texttt{ofAbelImage},
    \texttt{ofAbelLifts}, \texttt{ofChartsOfAbelImage},
    \texttt{ofChartsOfAbelLifts}, \texttt{ofPic0ClassSurjective}, and
    \texttt{PicRepDatum.toJacobianDataOfAbelLifts} are downstream packaging
    devices.  Likewise \texttt{Pic0RepresentabilityOverlap} and
    \texttt{Pic0RepresentabilityDescentData} start from an already supplied
    representation.  None replaces the quotient-carrier handoff in R8--R9.
  \item \texttt{blueprint/src/chapters/detail},
    \texttt{blueprint/src/chapters-legacy}, and
    \texttt{content-formalization.tex} are historical audit material.  They
    are not included by \texttt{content.tex}; any conflicting route, status,
    or aspirational proof in them is non-normative and must not be cited as
    project state.
  \end{enumerate}
\end{remark}

\section{Public-root and Challenge status}
\label{route-app-challenge-status}

The modules defining \texttt{finiteStageBaseChangeOverIso},
\texttt{pic0RepresentableBy\_finiteStageBaseChange},
\texttt{pic0ClassYoneda},
\texttt{pic0SepClosedClassYoneda\_app}, and
\texttt{pic0FiniteStageOrbitsInAffineOpen\_of\_isProjective} are present in
the source tree but not reachable from the public \texttt{AlgebraicJacobian}
root at this checkpoint.  They receive no \texttt{\textbackslash leanok}
credit in the normative chapters.  The first two remain only slice/post-base-
change comparisons, and the last remains conditional on projectivity, even after
future root integration.

The headline declarations \texttt{Jacobian}, \texttt{instGrpObj},
\texttt{smoothOfRelativeDimension\_genus}, the anonymous proper and
geometrically-irreducible instances, \texttt{ofCurve},
\texttt{comp\_ofCurve}, \texttt{exists\_unique\_ofCurve\_comp}, the map and
laws of \texttt{functor}, \texttt{baseChangeIso},
\texttt{baseChangeIso\_id}, \texttt{baseChangeIso\_comp}, and
\texttt{baseChange\_ofCurve} in \texttt{Challenge.lean} remain backed by
\texttt{sorry}.  \texttt{Jacobian.congr} is definitionally derived from the
sorry-bearing functor and is therefore axiom-contaminated rather than an
independent endpoint.  The generic scheme base-change functor and its
identity/composition isomorphisms are checked, but they do not prove the
Jacobian-specific comparison or coherences.

\section{Residual blocker ledger}

\begin{remark}[Blockers in dependency order]
  \label{route-app-blockers}
  \dcref{custom}

  \begin{enumerate}
  \item \textbf{B1, first unmet producer:} after taking the finite normal
    closure of all coefficients, finalize a legacy \(P\) with
    \(L=P.N.1/K\) finite Galois, construct \(\xi_L\) on its literal
    \texttt{P.gluedOver}, and prove (UC), including the exact big-etale and
    degree-zero transports.  Current tensor-stage work supplies local inputs
    only.
  \item \textbf{B2:} globalize affine faithful-flat Picard reflection to
    arbitrary tests.  After B1, use (UC) to prove double-overlap compatibility
    of representing morphisms, apply fpqc morphism descent, and reflect
    equality of classes.  This gives the all-test natural bijection (FS-Rep).
  \item \textbf{B3:} formalize the field-general algebraic-group
    quasi-projectivity theorem (0BF7) with its source prerequisites.
    After B2, apply it to the finite-type group on the exact \(J_L\),
    obtaining (QP).
  \item \textbf{B4:} apply homogeneous prime avoidance to (QP), then
    instantiate the finite-in-affine-to-orbits consumer to obtain (OA)
    for the canonical action on the same \(J_L\).
  \item \textbf{B5:} instantiate the conditional finite-Galois theorem and
    same-carrier handoff, producing one original-field
    \texttt{JacobianData C}.
  \item \textbf{B6:} construct the additive, scalar-compatible
    epsilon-kernel--\(H^1(C,\mathcal O_C)\) comparison (SL); the current
    set-level equivalence cannot prove the tangent dimension formula.
  \item \textbf{B7:} prove \texttt{abelCrossBaseCechCore} for every field
    extension and rational base point; the current graph theorem only reduces
    this to two Cech-class equations.
  \item \textbf{B8:} construct the fixed proper degree-\(n_C\) Abel source,
    geometric reducedness, and the relative-dimension-\(g\) zero-hypercover
    required by the checked geometry consumers.
  \item \textbf{B9:} prove the pointed genus-zero identification with
    \(\mathbb P^1\), all-test degree-zero Picard vanishing, and constancy/Hex;
    combine these with the positive-genus factorization to obtain (Alb).
  \item \textbf{B10:} select the datum family in \texttt{Challenge.lean}, prove
    functoriality and the Jacobian-specific base-change identity/tower
    coherences on that family, instantiate (AJ-BC), and expose all geometry
    without changing carriers or signatures.
  \end{enumerate}

  The arbitrary-test reflection lemma in B2 and the general group theorem
  in B3 are independent prerequisites.  Their applications to the selected
  carrier follow B1 and B2, respectively; B4--B5 consume those applications.
  B6--B10 concern the subsequent datum and retain the dependencies stated
  in Chapter~\ref{route-ch-handoff}.
\end{remark}
